% Source: Weinberg GR, physical PDF pp. 364--371
%         (printed pp. 342--349).
% Coverage: Section 11.9, Gravitational Collapse; equations
%           (11.9.1)--(11.9.46).
% Figures/tables/footnotes: no figures or tables; reference callouts 32--39.
% Status: source-reviewed and compile-clean.
% Uncertainties: the curvature-sign conversion is documented at
%                Eqs. (11.9.9)--(11.9.12); none unresolved.

\subsection{Gravitational Collapse}\label{sec:11.9}

We saw in Sections~\ref{sec:11.3} and \ref{sec:11.4} that a cooling star of
mass greater than a few solar masses cannot reach equilibrium as either a
white dwarf or a neutron star. It may be that a massive star will always
eject enough matter by the time it reaches the end of its thermonuclear
evolution so that its mass drops below the Chandrasekhar or the
Oppenheimer--Volkoff limits. If not, then it will collapse.

A proper treatment of gravitational collapse would be prohibitively
complicated for this book. In order to get some feeling for what can happen
during collapse, we consider only the simplest
case,\hyperref[ch11-ref-32]{\textsuperscript{32}} the spherically symmetric
collapse of ``dust'' with negligible pressure. Since the dust particles are
acted on by purely gravitational forces, they fall freely, and we can use them
as the physical basis of a comoving coordinate system of the sort discussed
in the last section. The metric is then given by Eq.~\eqref{eq:11.8.8}:
\begin{equation}
  \dd s^2
  =
  -\dd t^2+U(r,t)\,\dd r^2
  +V(r,t)(\dd\theta^2+\sin^2\theta\,\dd\varphi^2).
  \tag{11.9.1}\label{eq:11.9.1}
\end{equation}
The energy--momentum tensor for a fluid of negligible pressure is given by
Eq.~\eqref{eq:5.4.2} as
\begin{equation}
  T^{\mu\nu}=\rho U^\mu U^\nu,
  \tag{11.9.2}\label{eq:11.9.2}
\end{equation}
where \(\rho(r,t)\) is the proper energy density and \(U^\mu\) is the velocity
four-vector, given for a comoving coordinate system by
Eqs.~\eqref{eq:11.8.9} and \eqref{eq:11.8.11}:
\begin{equation}
  U^r=U^\theta=U^\varphi=0,\qquad U^t=1.
  \tag{11.9.3}\label{eq:11.9.3}
\end{equation}
The equations of momentum conservation
\(\nabla_\mu T_i{}^\mu=0\) are automatically satisfied, and the equation for
energy conservation reads
\[
  0=\nabla_\mu T_t{}^\mu
  =
  -\partial_t\rho-\rho\Gamma^\lambda{}_{\lambda t}
  =
  -\partial_t\rho
  -\rho\left(\frac{\dot U}{2U}+\frac{\dot V}{V}\right),
\]
or, in other words,
\begin{equation}
  \partial_t\!\left(\rho V\sqrt U\right)=0.
  \tag{11.9.4}\label{eq:11.9.4}
\end{equation}

In the binding curvature convention, the Einstein field equations can be
written
\begin{equation}
  R_{\mu\nu}=8\pi G S_{\mu\nu},
  \tag{11.9.5}\label{eq:11.9.5}
\end{equation}
where
\begin{equation}
  S_{\mu\nu}
  \equiv
  T_{\mu\nu}-\frac12g_{\mu\nu}T^\lambda{}_\lambda
  =
  \rho\left(\frac12g_{\mu\nu}+U_\mu U_\nu\right).
  \tag{11.9.6}\label{eq:11.9.6}
\end{equation}
This may be evaluated with the aid of Eqs.~\eqref{eq:11.9.1} and
\eqref{eq:11.9.3}; we find that the only nonvanishing components of
\(S_{\mu\nu}\) are
\begin{equation}
  S_{rr}=\frac{\rho U}{2},\qquad
  S_{\theta\theta}=\frac{\rho V}{2},\qquad
  S_{\varphi\varphi}=S_{\theta\theta}\sin^2\theta,\qquad
  S_{tt}=\frac{\rho}{2}.
  \tag{11.9.7}\label{eq:11.9.7}
\end{equation}
In particular,
\begin{equation}
  S_{tr}=0.
  \tag{11.9.8}\label{eq:11.9.8}
\end{equation}
Using Eqs.~\eqref{eq:11.9.7}--\eqref{eq:11.9.8} and
\eqref{eq:11.8.14}--\eqref{eq:11.8.17} in Eq.~\eqref{eq:11.9.5} yields
four field equations:
% MODERNIZATION CHECK: (11.9.9)--(11.9.12) display the target Ricci
% components, so both their left- and right-hand sides are the negatives
% of Weinberg's printed equations where appropriate.
\begin{align}
  &-\frac1U
  \left[
    \frac{V''}{V}
    -\frac{V'^2}{2V^2}
    -\frac{U'V'}{2UV}
  \right]
  +\frac{\ddot U}{2U}
  -\frac{\dot U^2}{4U^2}
  +\frac{\dot U\dot V}{2UV}
  =4\pi G\rho,
  \tag{11.9.9}\label{eq:11.9.9}\\
  &\frac1V-\frac1U
  \left[
    \frac{V''}{2V}
    -\frac{U'V'}{4UV}
  \right]
  +\frac{\ddot V}{2V}
  +\frac{\dot V\dot U}{4VU}
  =4\pi G\rho,
  \tag{11.9.10}\label{eq:11.9.10}\\
  &-\frac{\ddot U}{2U}
  -\frac{\ddot V}{V}
  +\frac{\dot U^2}{4U^2}
  +\frac{\dot V^2}{2V^2}
  =4\pi G\rho,
  \tag{11.9.11}\label{eq:11.9.11}\\
  &-\frac{\dot V'}{V}
  +\frac{V'\dot V}{2V^2}
  +\frac{\dot U V'}{2UV}
  =0.
  \tag{11.9.12}\label{eq:11.9.12}
\end{align}

Let us simplify our model even further, and assume that \(\rho\) is
independent of position.\hyperref[ch11-ref-32]{\textsuperscript{32}} We can
now seek a separable solution, with
\[
  U=a^2(t)f(r),\qquad V=b^2(t)g(r).
\]
Then Eq.~\eqref{eq:11.9.12} requires that \(\dot b/b\) equal \(\dot a/a\),
so we can normalize \(f\) and \(g\) so that
\[
  b(t)=a(t).
\]
Also, we are still free to redefine the radial coordinate as an arbitrary
function \(\tilde r\) of \(r\), and in particular we can choose
\(\tilde r=\sqrt{g(r)}\), so \(f\) and \(g\) are replaced with
\(\tilde f=fg'^2/(4g)\) and \(\tilde g=\tilde r^2\). Dropping the tildes,
we have then
\begin{equation}
  U=a^2(t)f(r),\qquad V=a^2(t)r^2.
  \tag{11.9.13}\label{eq:11.9.13}
\end{equation}
Equations~\eqref{eq:11.9.9} and \eqref{eq:11.9.10}, with both sides
multiplied by \(-a^2\), then read
\begin{align}
  -\frac{f'(r)}{rf^2(r)}
  -\ddot a(t)a(t)-2\dot a^2(t)
  &=-4\pi G a^2(t)\rho(t),
  \tag{11.9.14}\label{eq:11.9.14}\\
  \left[
    -\frac1{r^2}
    +\frac1{r^2f(r)}
    -\frac{f'(r)}{2rf^2(r)}
  \right]
  -\ddot a(t)a(t)-2\dot a^2(t)
  &=-4\pi G a^2(t)\rho(t).
  \tag{11.9.15}\label{eq:11.9.15}
\end{align}
The first terms in Eqs.~\eqref{eq:11.9.14} and \eqref{eq:11.9.15} must
evidently be equal constants, which we shall call \(-2k\):
\[
  -\frac{f'(r)}{rf^2(r)}
  =
  -\frac1{r^2}+\frac1{r^2f(r)}
  -\frac{f'(r)}{2rf^2(r)}
  =-2k.
\]
The unique solution is
\[
  f(r)=(1-kr^2)^{-1},
\]
so the metric takes the form
\begin{equation}
  \dd s^2
  =
  -\dd t^2+a^2(t)
  \left[
    \frac{\dd r^2}{1-kr^2}
    +r^2\dd\theta^2+r^2\sin^2\theta\,\dd\varphi^2
  \right].
  \tag{11.9.16}\label{eq:11.9.16}
\end{equation}
(Incidentally, this metric is spatially homogeneous as well as isotropic, and
for this reason it will provide the kinematic framework for our treatment of
relativistic cosmology in Chapter~14.)

Our remaining problem is to calculate the functions \(\rho(t)\) and \(a(t)\).
Using Eqs.~\eqref{eq:11.9.13} and \eqref{eq:11.9.14} in the
energy-conservation equation \eqref{eq:11.9.4}, we find that
\(\rho(t)a^3(t)\) is constant. We normalize the radial coordinate \(r\) so
that
\begin{equation}
  a(0)=1,
  \tag{11.9.17}\label{eq:11.9.17}
\end{equation}
and therefore
\begin{equation}
  \rho(t)=\rho(0)a^{-3}(t).
  \tag{11.9.18}\label{eq:11.9.18}
\end{equation}
The field equations \eqref{eq:11.9.14} or \eqref{eq:11.9.15} and
\eqref{eq:11.9.11} are now ordinary differential equations:
\begin{align}
  -2k-\ddot a(t)a(t)-2\dot a^2(t)
  &=-4\pi G\rho(0)a^{-1}(t),
  \tag{11.9.19}\label{eq:11.9.19}\\
  \ddot a(t)a(t)
  &=-\frac{4\pi G}{3}\rho(0)a^{-1}(t).
  \tag{11.9.20}\label{eq:11.9.20}
\end{align}
We can eliminate \(\ddot a(t)\) by adding these two equations, and find
\begin{equation}
  \dot a^2(t)
  =
  -k+\frac{8\pi G}{3}\rho(0)a^{-1}(t).
  \tag{11.9.21}\label{eq:11.9.21}
\end{equation}
Equations~\eqref{eq:11.9.19} and \eqref{eq:11.9.20} can be recovered from
Eq.~\eqref{eq:11.9.21} and its time derivative, so we can forget about them
and simply use Eq.~\eqref{eq:11.9.21} to calculate \(a(t)\).

We shall now assume that the fluid is at rest (in standard coordinates) at
\(t=0\), so
\begin{equation}
  \dot a(0)=0,
  \tag{11.9.22}\label{eq:11.9.22}
\end{equation}
and therefore Eqs.~\eqref{eq:11.9.21} and \eqref{eq:11.9.17} give
\begin{equation}
  k=\frac{8\pi G}{3}\rho(0).
  \tag{11.9.23}\label{eq:11.9.23}
\end{equation}
Thus Eq.~\eqref{eq:11.9.21} can be written
\begin{equation}
  \dot a^2(t)=k[a^{-1}(t)-1].
  \tag{11.9.24}\label{eq:11.9.24}
\end{equation}
The solution is given by the parametric equations of a cycloid:
\begin{equation}
  t=\frac{\psi+\sin\psi}{2\sqrt k},
  \qquad
  a=\frac12(1+\cos\psi).
  \tag{11.9.25}\label{eq:11.9.25}
\end{equation}
Note that \(a(t)\) vanishes when \(\psi=\pi\), and hence when \(t=T\), where
\begin{equation}
  T=\frac{\pi}{2\sqrt k}
  =\frac{\pi}{2}
   \left[\frac{3}{8\pi G\rho(0)}\right]^{1/2}.
  \tag{11.9.26}\label{eq:11.9.26}
\end{equation}
\emph{Thus a fluid sphere of initial density \(\rho(0)\) and zero pressure
will collapse from rest to a state of infinite proper energy density in the
finite time \(T\).}

Although the collapse is complete at a finite coordinate time \(t=T\), any
light signal coming to us from the sphere's surface will be delayed by its
gravitational field (see Section~\ref{sec:8.7}), so we on earth will not see
the star suddenly vanish. To make this more specific, we have to complete our
calculation by finding the metric outside the star.

The Birkhoff theorem proved in Section~\ref{sec:11.7} shows that it is always
possible to find a ``standard'' coordinate system
\(\bar r,\bar\theta,\bar\varphi,\bar t\) in which the metric outside the
sphere takes the form
\begin{equation}
  \dd s^2
  =
  -\left(1-\frac{2MG}{\bar r}\right)\dd\bar t^2
  +\left(1-\frac{2MG}{\bar r}\right)^{-1}\dd\bar r^2
  +\bar r^2(\dd\bar\theta^2+\sin^2\bar\theta\,\dd\bar\varphi^2).
  \tag{11.9.27}\label{eq:11.9.27}
\end{equation}
But this metric is not in the Gaussian normal form \eqref{eq:11.9.1}, so in
order to match solutions at the surface we either have to convert the
interior solution \eqref{eq:11.9.16} into standard coordinates, or the
exterior solution \eqref{eq:11.9.27} into Gaussian normal coordinates. We
choose the former course.\hyperref[ch11-ref-32]{\textsuperscript{32}}

Inspection of Eq.~\eqref{eq:11.9.16} shows immediately that the standard
spatial coordinates \(\bar r,\bar\theta,\bar\varphi\) must be chosen as
\begin{equation}
  \bar r=ra(t),\qquad
  \bar\theta=\theta,\qquad
  \bar\varphi=\varphi.
  \tag{11.9.28}\label{eq:11.9.28}
\end{equation}
In order to define a standard time coordinate such that \(\dd s^2\) does not
contain a cross-term \(\dd\bar r\,\dd\bar t\), we employ the ``integrating
factor'' technique described in Section~\ref{sec:11.7}, which gives
\begin{equation}
  \bar t
  =
  \left(\frac{1-kr_\star^2}{k}\right)^{1/2}
  \int_{S(r,t)}^1
  \frac{\dd\alpha}{1-kr_\star^2/\alpha}
  \left(\frac{\alpha}{1-\alpha}\right)^{1/2},
  \tag{11.9.29}\label{eq:11.9.29}
\end{equation}
where
\begin{equation}
  S(r,t)
  =
  1-
  \left(\frac{1-kr^2}{1-kr_\star^2}\right)^{1/2}
  [1-a(t)].
  \tag{11.9.30}\label{eq:11.9.30}
\end{equation}
The constant \(r_\star\) is arbitrary, but may conveniently be chosen as the
radius of the sphere in comoving coordinates. It is straightforward to check
that the metric in the coordinate system
\(\bar r,\bar\theta,\bar\varphi,\bar t\) takes the standard form
\[
  \dd s^2
  =
  -B(\bar r,\bar t)\,\dd\bar t^2
  +A(\bar r,\bar t)\,\dd\bar r^2
  +\bar r^2(\dd\bar\theta^2+\sin^2\bar\theta\,\dd\bar\varphi^2),
\]
with
\begin{align}
  B
  &=
  \frac{a(t)}{S}
  \left(\frac{1-kr^2}{1-kr_\star^2}\right)^{1/2}
  \frac{(1-kr_\star^2/S)^2}{1-kr^2/a(t)},
  \tag{11.9.31}\label{eq:11.9.31}\\
  A
  &=
  \left(1-\frac{kr^2}{a(t)}\right)^{-1}.
  \tag{11.9.32}\label{eq:11.9.32}
\end{align}
It is now understood that \(S\) is a function of \(\bar t\) defined by
Eq.~\eqref{eq:11.9.29}, and that \(r\) and \(a(t)\) are functions of
\(\bar r\) and \(S\), or \(\bar r\) and \(\bar t\), defined by solving
Eqs.~\eqref{eq:11.9.28} and \eqref{eq:11.9.30}. This is a mess, but at the
radius \(r=r_\star\) of the star (a constant, since \(r\) is a comoving
coordinate) we have
\begin{align}
  \bar r&=\bar r_\star(t)\equiv r_\star a(t),
  \tag{11.9.33}\label{eq:11.9.33}\\
  \bar t
  &=
  \left(\frac{1-kr_\star^2}{k}\right)^{1/2}
  \int_{a(t)}^1
  \frac{\dd\alpha}{1-kr_\star^2/\alpha}
  \left(\frac{\alpha}{1-\alpha}\right)^{1/2},
  \tag{11.9.34}\label{eq:11.9.34}\\
  B(\bar r_\star,\bar t)&=1-\frac{kr_\star^2}{a(t)},
  \tag{11.9.35}\label{eq:11.9.35}\\
  A(\bar r_\star,\bar t)&=
  \left(1-\frac{kr_\star^2}{a(t)}\right)^{-1}.
  \tag{11.9.36}\label{eq:11.9.36}
\end{align}
[Equation~\eqref{eq:11.9.34} could have been obtained by integrating the
equations for free fall given in Section~\ref{sec:8.4}.] Comparing with
Eq.~\eqref{eq:11.9.27}, we see that the interior and exterior solutions fit
continuously at \(\bar r=r_\star a(t)\) if
\begin{equation}
  k=\frac{2MG}{r_\star^3}.
  \tag{11.9.37}\label{eq:11.9.37}
\end{equation}
With Eq.~\eqref{eq:11.9.23}, this just says that
\begin{equation}
  M=\frac{4\pi}{3}\rho(0)r_\star^3,
  \tag{11.9.38}\label{eq:11.9.38}
\end{equation}
not a surprising result!

Now we return to the problem of calculating the behavior of light signals
emitted from the surface of the collapsing sphere. A light signal emitted in
a radial direction at a standard time \(\bar t\) will have
\(\dd\bar r/\dd\bar t\) given by Eq.~\eqref{eq:11.9.27} and the condition
\(\dd s^2=0\), so it will arrive at a distant point \(\bar r'\) at a time
\begin{equation}
  \bar t'
  =
  \bar t+
  \int_{r_\star a(t)}^{\bar r'}
  \left(1-\frac{2MG}{\bar r}\right)^{-1}\dd\bar r.
  \tag{11.9.39}\label{eq:11.9.39}
\end{equation}
The most striking consequence of Eqs.~\eqref{eq:11.9.39} and
\eqref{eq:11.9.34} is that both \(\bar t\) and \(\bar t'\) approach infinity
when the radius \eqref{eq:11.9.33} of the sphere approaches the Schwarzschild
radius \(2GM\), that is, when
\begin{equation}
  a(t)\longrightarrow\frac{2GM}{r_\star}=kr_\star^2.
  \tag{11.9.40}\label{eq:11.9.40}
\end{equation}
\emph{The collapse to the Schwarzschild radius therefore appears to an
outside observer to take an infinite time, and the collapse to \(a=0\) is
utterly unobservable from outside.}

Although the collapsing sphere does not suddenly disappear, it does fade out
of sight, because light from its surface is subject to an increasing redshift.
The proper time for a light source on the sphere's surface is just the
comoving time \(t\), so the comoving time interval between emission of wave
crests at the surface equals the natural wavelength \(\lambda_0\) that would
be emitted by the source in the absence of gravitation. The standard time
interval \(\dd\bar t'\) between arrivals of wave crests at \(\bar r'\) is the
observed wavelength \(\lambda'\); thus the fractional change of wavelength is
\[
  \begin{aligned}
  z
  &\equiv\frac{\lambda'-\lambda_0}{\lambda_0}
   =\frac{\dd\bar t'}{\dd t}-1\\
  &=\frac{\dd\bar t}{\dd t}
    -r_\star\dot a(t)
     \left(1-\frac{2MG}{r_\star a(t)}\right)^{-1}
    -1\\
  &=-\dot a(t)
    \left(1-\frac{kr_\star^2}{a(t)}\right)^{-1}
    \left[
      \left(\frac{1-kr_\star^2}{k}\right)^{1/2}
      \left(\frac{a(t)}{1-a(t)}\right)^{1/2}
      +r_\star
    \right]-1.
  \end{aligned}
\]
Using Eq.~\eqref{eq:11.9.24} to determine \(\dot a(t)\), this is
\begin{equation}
  z
  =
  \left(1-\frac{kr_\star^2}{a(t)}\right)^{-1}
  \left[
    (1-kr_\star^2)^{1/2}
    +r_\star\sqrt k
     \left(\frac{1-a(t)}{a(t)}\right)^{1/2}
  \right]-1.
  \tag{11.9.41}\label{eq:11.9.41}
\end{equation}
In order to see how the redshift \(z\) varies with \(\bar t'\), let us assume
that the sphere is initially very much larger than its Schwarzschild radius,
\begin{equation}
  kr_\star^2=\frac{2GM}{r_\star}\ll1,
  \tag{11.9.42}\label{eq:11.9.42}
\end{equation}
and distinguish two periods in the history of the collapse:

\emph{(A)} Until \(t\) gets close to \(T\), we have
\begin{equation}
  \frac{kr_\star^2}{a(t)}\ll1.
  \tag{11.9.43}\label{eq:11.9.43}
\end{equation}
Using Eqs.~\eqref{eq:11.9.42} and \eqref{eq:11.9.43} in
\eqref{eq:11.9.34}, \eqref{eq:11.9.39}, and \eqref{eq:11.9.41} gives
(with \(\bar r'\gg r_\star\))
\begin{equation}
  \begin{aligned}
  \bar t&\simeq t,\\
  \bar t'
  &\simeq\bar t+\bar r'-r_\star a(t)
   \simeq t+\bar r'-r_\star a(t)
   \simeq t+\bar r',\\
  z
  &\simeq
  r_\star\sqrt k
  \left(\frac{1-a(t)}{a(t)}\right)^{1/2}
  \simeq
  r_\star\sqrt k
  \left[
    \frac{1-a(\bar t'-\bar r')}{a(\bar t'-\bar r')}
  \right]^{1/2}.
  \end{aligned}
  \tag{11.9.44}\label{eq:11.9.44}
\end{equation}

\emph{(B)} Eventually we have
\[
  \frac{kr_\star^2}{a(t)}\longrightarrow1
\]
at a time \(t_1\) given by Eq.~\eqref{eq:11.9.25} as
\begin{equation}
  t_1
  \simeq
  \frac1{2\sqrt k}
  \left[\pi-\frac43(kr_\star^2)^{3/2}\right].
  \tag{11.9.45}\label{eq:11.9.45}
\end{equation}
Now Eqs.~\eqref{eq:11.9.34}, \eqref{eq:11.9.39}, and
\eqref{eq:11.9.41} give
\begin{equation}
  \begin{aligned}
  \bar t
  &\simeq
  -kr_\star^3\ln\left[1-\frac{kr_\star^2}{a(t)}\right]
  +\text{constant},\\
  \bar t'
  &\simeq
  \bar t-kr_\star^3\ln\left[1-\frac{kr_\star^2}{a(t)}\right]
  +\text{constant}\\
  &\simeq
  -2kr_\star^3\ln\left[1-\frac{kr_\star^2}{a(t)}\right]
  +\text{constant},\\
  z
  &\simeq
  2\left(1-\frac{kr_\star^2}{a(t)}\right)^{-1}
  \propto
  \exp\left(\frac{\bar t'}{2kr_\star^3}\right).
  \end{aligned}
  \tag{11.9.46}\label{eq:11.9.46}
\end{equation}
Putting (A) and (B) together, we conclude that the redshift \(z\) seen by an
observer at \(\bar r'\) is zero when the collapse is observed to begin,
increases gradually but remains of order \(r_\star\sqrt k\ll1\) until a time
very
close to \(T=\pi/(2\sqrt k)\) has passed, and then grows exponentially with
a rate \(1/(2kr_\star^3)\). For example, a collapsing sphere with a mass
\(M=10^8M_\odot\) and radius \(r_\star=100\) light-years will have a redshift
\(z\) of order \(10^{-3}\) for a period of order \(10^5\) years, after which
the redshift suddenly begins growing exponentially with an e-folding time of
order \(1\) minute. For practical purposes, the collapsing sphere is suddenly
and completely cut off from communication with the rest of the universe.

Completely cut off? Even if a collapsing body does fade out of sight, it
still has a gravitational field, and, as shown in Section~\ref{sec:7.6}, the
measurement of this field at great distances can be used to determine the
energy, momentum, and angular momentum of the body. If the body has a net
electric charge, then measurement of the electric field at great distances
will, via Gauss's theorem, also tell us the charge. It is interesting to ask
whether measurements of the gravitational and/or electromagnetic fields
outside a collapsing body can yield any information about the body beyond
the energy, momentum, angular momentum, and charge. In the case of a
spherically symmetric electrically neutral body, which we have been
considering in this chapter, the answer is provided by Birkhoff's theorem:
the gravitational field outside a spherically symmetric body must be of the
Schwarzschild form, so all we can ever learn about the body is its mass.
(Spherical symmetry, of course, implies zero momentum and zero angular
momentum.)

Carter\hyperref[ch11-ref-33]{\textsuperscript{33}} has shown that, when the
gravitational field of an \emph{axially symmetric} collapsing body settles
down to a stationary state, its exterior metric belongs to a
\emph{two-parameter} family of solutions, such as the Kerr metrics (see
Section~\ref{sec:11.7}) that are completely specified by the total mass and
angular momentum. It is widely believed that the gravitational field of any
electrically neutral collapsing body will eventually approach the Kerr form.

As mentioned in the introduction to this chapter, interest in the phenomenon
of gravitational collapse was rekindled in the last decade by the discovery
of quasi-stellar sources, which appear to require some powerful new source of
energy. The maximum energy available from fusion of hydrogen into the most
stable nuclei, say iron, is only \(8\,\mathrm{MeV}\) per nucleon, or less than
\(1\%\) of the rest mass. Matter--antimatter annihilation could have
\(100\%\) efficiency (apart from neutrino energy losses), but this process can
be important only if there is some abundant natural source of antinucleons.
Otherwise the only likely mechanism for conversion of mass into energy with
high efficiency is gravitational
collapse.\hyperref[ch11-ref-34]{\textsuperscript{34}}

A cloud of dust that is collapsing as in the Oppenheimer--Snyder model will
obviously release no energy to the outside world. To extract the growing
kinetic energy of the falling particles, something must slow them on the way
down---either a macroscopic ``bounce'' of the whole system, or
particle--particle collisions that heat the collapsing gas. Detailed
calculations reveal a discouragingly low efficiency for conversion of mass
into available energy in the gravitational collapse of an \emph{isolated}
body.\hyperref[ch11-ref-35]{\textsuperscript{35}} However, particles falling
into a Kerr metric can reemerge with a higher energy, acquired at the expense
of the rotational energy of the collapsing
body.\hyperref[ch11-ref-36]{\textsuperscript{36}}

Whether or not gravitational collapse has anything to do with quasi-stellar
sources, the question remains: What happens to a real cooling star whose mass
is above the Chandrasekhar and Oppenheimer--Volkoff limits? In recent years
topological methods have been used by Penrose and Hawking to prove a number
of powerful theorems,\hyperref[ch11-ref-37]{\textsuperscript{37}} to the
effect that under reasonable conditions (validity of general relativity,
positivity of energy, ubiquity of matter, causality) collapse becomes
inevitable once a \emph{trapped surface} forms. A trapped surface is a closed
spacelike two-dimensional surface for which both the outgoing and the ingoing
families of future-directed null geodesics orthogonal to the surface are
converging. (For the Schwarzschild metric, the spheres with \(r\) and \(t\)
constant are trapped surfaces for \(r\) within the Schwarzschild radius
\(2MG\).) However, it is not known whether a real massive star will actually
develop a trapped surface, or merely explode into fragments with small enough
mass to form stable neutron stars or white dwarfs.

If gravitational collapse is indeed the inevitable fate of massive bodies,
then we must expect that the universe is full of \emph{black holes},
collapsing bodies whose presence is betrayed only through their gravitational
fields or through the energy released when matter is drawn
in.\hyperref[ch11-ref-38]{\textsuperscript{38}} Our best hope of observing
gravitational collapse would be to find a binary star, one member an ordinary
visible star, and the other member a black
hole.\hyperref[ch11-ref-39]{\textsuperscript{39}}
